%%%%%%%%%%%%%%%%%%%%%%%%%%%%%%%%%%%%%%%%%%%%%%%%%%%%%%%%%%%%%%%%%%%%%%%%%%%%
% Manuscript targeted at the Journal of Statistics and Data Science Education
% (JSDSE; American Statistical Association / Taylor & Francis).
%
% SUBMISSION NOTES
% ----------------
% * JSDSE uses DOUBLE-ANONYMOUS peer review (confirmed via the journal's
%   DOAJ record). The \anon switch below controls whether the manuscript
%   renders in anonymized (review) or identified (final) form. Compile with
%   \anontrue for anything you upload to ScholarOne for review.
% * Submission is through ScholarOne: https://mc.manuscriptcentral.com/ujse
% * Taylor & Francis operates format-free submission for most titles, so a
%   clean, consistently-cited, double-spaced manuscript in the standard
%   article class (as here) is acceptable at initial submission. The
%   official T&F class is only needed if the editors request formatting at
%   revision or acceptance.
% * TO SWITCH TO THE OFFICIAL T&F CLASS LATER: replace the \documentclass
%   line below with \documentclass[largeformat]{interact} and switch the
%   bibliography style to the T&F author-year style shipped with the
%   template. NOTE: interact.cls is NOT on CTAN and is not part of TeX
%   Live -- it must be downloaded from Taylor & Francis Author Services, or
%   obtained by opening the "Taylor & Francis (Interact layout + APA
%   reference style)" template on Overleaf.
% * This is NOT based on overleaf-template/main.tex: that template is built
%   for accounting and finance journals (12pt double-spaced article,
%   biblatex-chicago, hypothesis environments) and its house style does not
%   match ASA/T&F conventions.
%%%%%%%%%%%%%%%%%%%%%%%%%%%%%%%%%%%%%%%%%%%%%%%%%%%%%%%%%%%%%%%%%%%%%%%%%%%%

\documentclass[12pt]{article}

%%%%%%%%%%%%%%%%%%%%%%%%%%%%%%%%%%%%%%%%%%%%%%%%%%%%%%%%%%%%%%%%%%%%%%%%%%%%
% ANONYMIZATION SWITCH
% \anontrue  -> double-anonymous review version (author, institution,
%               repository URLs, DOI, and acknowledgements masked)
% \anonfalse -> identified version (for the accepted/published manuscript,
%               for a preprint, or for your own records)
%%%%%%%%%%%%%%%%%%%%%%%%%%%%%%%%%%%%%%%%%%%%%%%%%%%%%%%%%%%%%%%%%%%%%%%%%%%%
\newif\ifanon
\anontrue

\usepackage[margin=1in]{geometry}
\usepackage{lmodern}      % Type 1 outline fonts (avoids bitmap-font PDFs)
\usepackage[T1]{fontenc}
\usepackage[utf8]{inputenc}
\usepackage{setspace}
\usepackage{booktabs}
\usepackage{tabularx}
\usepackage{graphicx}
\usepackage{amsmath}
\usepackage[round,authoryear]{natbib}
\usepackage{enumitem}
\usepackage{url}
\usepackage[hidelinks]{hyperref}

\bibliographystyle{apalike}

\doublespacing
\setlength{\parindent}{0.5in}

% Masked stand-ins used only in the anonymized build.
\ifanon
  \newcommand{\hubsite}{[hub site URL omitted for anonymous review]}
  \newcommand{\pipelinerepo}{[pipeline template URL omitted for anonymous review]}
  \newcommand{\manuscriptrepo}{[manuscript template URL omitted for anonymous review]}
  \newcommand{\institution}{a large public research university in the United States}
\else
  \newcommand{\hubsite}{\url{https://eweisbrod.github.io/example-project/}}
  \newcommand{\pipelinerepo}{\url{https://github.com/eweisbrod/project-template}}
  \newcommand{\manuscriptrepo}{\url{https://github.com/eweisbrod/overleaf-template}}
  \newcommand{\institution}{the University of Kansas}
\fi

\title{\bfseries Teaching Reproducible Empirical Research Workflows in
Accounting and Finance: An Open-Source Template Suite}

\ifanon
  \author{}
  \date{}
\else
  \author{Eric H. Weisbrod\thanks{School of Business, University of Kansas.
  Email: \href{mailto:eric.weisbrod@ku.edu}{eric.weisbrod@ku.edu}.
  ORCID: 0000-0002-8814-250X.}}
  \date{August 2026}
\fi

\begin{document}

\maketitle

\begin{abstract}
\noindent
Doctoral students in accounting and finance are increasingly expected to
produce research that satisfies formal data and code sharing policies at
the field's major journals, yet few doctoral programs offer formal
training in the end-to-end research workflow those policies presuppose.
This article describes an open-source suite of teaching materials and
repository templates for reproducible empirical research, developed for
and refined through doctoral teaching. The materials implement a complete
five-step research pipeline---licensed-database download, data
transformation, figures, analysis tables, and provenance
reporting---in parallel R, Python, and Stata implementations, and produce
the artifact set required by modern journal policies (execution logs,
sample-identifier files, and content-addressed file inventories) as a
by-product of running the pipeline. We describe the instructional design
of the materials, their use in a doctoral seminar, and the distinctive
pedagogical constraints of teaching reproducibility in a field whose raw
data are licensed and cannot be redistributed. The materials are freely
available under a Creative Commons Attribution license and are designed
for adoption by instructors in business schools and adjacent
observational-data disciplines.

\vspace{1em}
\noindent\textbf{Keywords:} reproducibility; research workflow; version
control; doctoral education; licensed data; multi-language instruction
\end{abstract}

\newpage

\section{Introduction}

The statistics and data science education community has argued for more
than a decade that computing workflow deserves explicit instructional
attention rather than being left for students to absorb implicitly
\citep{nolan2010computing, cetinkayarundel2018infrastructure}. More
recently, reproducibility and ``responsible workflow'' have been advanced
as core curricular outcomes in their own right
\citep{horton2022growing}, with concrete proposals for teaching version
control \citep{bryan2018version, beckman2021git} and for making teaching
materials themselves reproducible \citep{dogucu2022reproducible}. This
journal has also formalized its own expectations for authors to share
deidentified data and code \citep{horton2024editorial}.

This article extends that agenda to an audience the education literature
rarely reaches: doctoral students and early-career researchers in
accounting and finance. Empirical research in these fields sits on a
workflow stack that will look familiar to statistics educators---database
access, data wrangling, statistical modeling, and publication-grade table
and figure production---but with three distinctive features that shape
how reproducible workflow must be taught.

First, the raw data are almost universally licensed. The standard
academic data provider, Wharton Research Data Services (WRDS), licenses
commercial databases (Compustat, CRSP, IBES) whose contents cannot
legally be redistributed, so the canonical reproducibility
advice---publish the data alongside the code
\citep{peng2011reproducible}---does not directly apply. Second, the field
is genuinely multi-language: research teams routinely mix SAS, Stata, R,
and Python within a single project, often splitting data preparation and
analysis across languages. Third, and most consequentially for students,
the field's journals have begun to formalize data and code sharing
policies modeled on norms in economics
\citep{christensen2018transparency}. The \textit{Journal of Accounting
Research}, for example, now requires submitting authors to provide the
code that produces the analytical dataset and all reported tables and
figures from raw inputs, a comprehensive log file documenting end-to-end
execution, and identifiers of the observations comprising the final
sample \citep{JARpolicy}. Comparable requirements have been adopted at
several other major journals in the field.

These policies create an authentic, high-stakes motivation for workflow
instruction: a doctoral student who cannot package a research pipeline to
policy standard cannot publish in the field's top journals. Yet surveys
of the accounting research community suggest that reproducibility
practices and training lag the field's stated aspirations
\citep{hail2020reproducibility}, and formal workflow training remains
uncommon in business doctoral programs. Most students arrive with no
software development background and assemble their workflow habits by
imitation---inheriting scripts, conventions, and bad habits from advisors
and coauthors.

The materials described here address this gap with a suite of
open-source teaching resources and repository templates that let a
student stand up a complete, policy-compliant empirical project and learn
the workflow conventions by working inside it. The suite takes the
position that workflow is best taught through a full-scale worked example
rather than through isolated exercises: students clone a working
five-step pipeline built on a real research design---a quarterly
earnings-announcement event study in the tradition of
\citet{ballbrown1968}---run it end to end against WRDS, and then extend
it. The materials were developed for, and have been refined through,
doctoral teaching at \institution.

The remainder of the article proceeds as follows. Section~\ref{sec:bg}
situates the materials in the reproducibility-education literature and
the disciplinary context. Section~\ref{sec:materials} describes the
architecture and content of the materials. Section~\ref{sec:design}
presents the learning objectives and the instructional design decisions
behind them. Section~\ref{sec:classroom} describes classroom use.
Section~\ref{sec:discussion} discusses adoptability, the pedagogy of
reproducibility under licensed data, and limitations.

\section{Background}
\label{sec:bg}

\subsection{Reproducibility and workflow in the education literature}

The case for teaching computing workflow explicitly, rather than assuming
students will acquire it osmotically, was made forcefully by
\citet{nolan2010computing}, who argued that computing is a core
statistical science competency rather than a service skill.
\citet{cetinkayarundel2018infrastructure} describe infrastructure for
integrating computing throughout a statistics curriculum, and
\citet{bryan2018version} and \citet{beckman2021git} make the specific
case for version control with Git and GitHub as a learning objective,
including concrete course designs at multiple levels.
\citet{dogucu2022reproducible} extend the argument from research to
teaching itself, recommending tools and practices for reproducible course
materials. The editorial position of this journal has described
reproducibility and responsible workflow as of ``growing importance''
across the data science and statistics curriculum
\citep{horton2022growing}. In parallel, the scientific computing
community has converged on pragmatic workflow guidance for researchers
who are not software engineers, most notably the ``good enough
practices'' of \citet{wilson2017good}: organize projects predictably,
script everything, separate raw from derived data, and put manual steps
under version control.

The materials described here operationalize that guidance for a specific
disciplinary audience. What they add to the existing literature is not
new workflow principles but a complete, domain-authentic implementation:
the gap between ``use version control and script your pipeline'' and a
working, policy-compliant empirical accounting project is large, and the
materials are designed to close it.

\subsection{The disciplinary context}

Two features of the accounting and finance research environment deserve
elaboration for readers outside these fields.

\paragraph{Journal artifact requirements.}
The field's flagship journals increasingly require a specific artifact
set at submission or acceptance. The \textit{Journal of Accounting
Research} policy \citep{JARpolicy} is representative and concrete: (1)
code that produces the analytical dataset, tables, and figures from raw
inputs; (2) a comprehensive log file documenting end-to-end execution;
and (3) identifiers of the observations in the final sample. The logic of
the policy reflects the licensed-data reality: because authors cannot
deposit Compustat or CRSP extracts, verifiability rests on the code, the
logs, and the sample identifiers, which allow any researcher with a WRDS
license to reconstruct the sample and reproduce the results. This
artifact set effectively defines a set of terminal workflow competencies
for doctoral students, and the materials described here are
backward-designed from it.

\paragraph{Multi-language practice.}
Legacy code in the field is predominantly SAS; Stata remains standard for
econometric estimation in many groups; R and Python adoption is growing
rapidly. A student entering a coauthored project has little control over
the language mix and benefits from structural
conventions---project layout, configuration, logging---that transcend any
one language. Recent field-specific pedagogical resources, notably
\citet{gow2024empirical}, teach empirical accounting research in R; the
materials here are complementary, focusing on the surrounding workflow
and deliberately spanning languages.

\section{The materials}
\label{sec:materials}

\subsection{Architecture: a hub and two templates}

The materials are organized as three public repositories that together
cover the workflow from raw database access to publication-grade
manuscript production:

\begin{itemize}[leftmargin=*]
  \item \textbf{The teaching hub} (\hubsite) contains the introductory
  landing content and seven chapter-length topic notes covering project
  structure, Git and GitHub for research, integrated development
  environment (IDE) setup, Python virtual environments, environment
  variables and configuration, AI-assistant integration, and a SAS macro
  reference for projects that retain SAS code.

  \item \textbf{The pipeline template} (\pipelinerepo) is a repository
  template implementing the five-step empirical pipeline with parallel
  implementations in R \citep{RCoreTeam} and Python, plus a Stata
  \citep{Stata} implementation of the analysis-and-tables step. At first
  run, a setup script prompts the user to select a language combination
  (full R, full Python, or mixed combinations including Stata) and prunes
  the unused files, so each student works in a clean single- or
  dual-language project rather than a cluttered polyglot one.

  \item \textbf{The manuscript template} (\manuscriptrepo) is a
  \LaTeX{} template designed to consume the table and figure outputs of
  the pipeline template, configured for real-time multi-author editing
  through a cloud \LaTeX{} platform, and demonstrating author-year
  citation handling, hypothesis-numbering conventions, and a section
  structure suitable for the field's journals.
\end{itemize}

The templates use the host platform's ``use this template'' workflow
rather than forking, so each student's project starts with a clean
history and no upstream link---mirroring how the template is meant to be
used for real research projects. All materials are released under the
Creative Commons Attribution 4.0 International (CC-BY-4.0) license, and
the hub is archived with a citable DOI.

\subsection{The running example and the five-step pipeline}

The running example throughout the materials is a quarterly
earnings-announcement event study: three-day buy-and-hold abnormal
returns around earnings announcements are regressed on unexpected
earnings interacted with a same-sign indicator on the seasonal change in
sales. The design descends from the foundational information-content
literature \citep{ballbrown1968}, so doctoral students in accounting and
finance recognize it immediately as ``real'' research rather than a toy;
at the same time it is compact enough to run end to end in minutes. The
example requires merging the two most important databases in the field
(Compustat fundamentals and CRSP security prices) through the
CRSP/Compustat linking table---a join that nearly every empirical project
in the field performs and nearly every student initially gets wrong.

The pipeline is organized as five numbered steps
(Table~\ref{tab:pipeline}), a structure that encodes several of the ``good
enough practices'' of \citet{wilson2017good} directly into the project
skeleton: raw and derived data are separated, every transformation is
scripted, and step numbers communicate execution order. The hub's topic
chapters (Table~\ref{tab:chapters}) cover the surrounding infrastructure
that the pipeline assumes.

\begin{table}[htbp]
\centering
\caption{Topic chapters in the teaching hub.}
\label{tab:chapters}
\small
\begin{tabularx}{\textwidth}{@{}lX@{}}
\toprule
\textbf{Chapter} & \textbf{Content} \\
\midrule
Project structure & Storage locations (repository host, local clone,
cloud sync, manuscript platform); separation of code, data, and
manuscript; folder layout and file-naming conventions. \\
Git and GitHub & Version-control workflow for research; commit, branch,
and tag conventions; ignore-file essentials; collaboration via pull
requests. \\
Setting up your IDE & Editor selection across RStudio, VS Code, Cursor,
and Positron; per-IDE configuration; AI-assistant integration. \\
Python virtual environments & Why Python requires environment isolation
while R, Stata, and SAS typically do not; the \texttt{uv} tool for
environment and dependency management. \\
Environment variables & Project-level configuration through a
\texttt{.env} file shared across R, Python, Stata, and SAS. \\
AI assistant context files & Conventions for providing project context to
AI coding assistants. \\
SAS macros & A reference macro library and a working example script
demonstrating integration with the templates' logging conventions. \\
\bottomrule
\end{tabularx}
\end{table}

\begin{table}[htbp]
\centering
\caption{Pipeline steps in the template, with language coverage.}
\label{tab:pipeline}
\small
\begin{tabularx}{\textwidth}{@{}lXl@{}}
\toprule
\textbf{Step} & \textbf{Content} & \textbf{Languages} \\
\midrule
001 Download data & Pulls Compustat fundamentals, CRSP returns, and
linking tables from WRDS; streams large tables via server-side cursors to
bounded-memory parquet files. & R, Python \\
002 Transform data & Constructs analytical variables, applies sample
filters, and computes event-window buy-and-hold abnormal returns. & R,
Python \\
003 Figures & Produces five publication-grade figures via
\texttt{ggplot2} \citep{ggplot2} in R and \texttt{plotnine}
\citep{plotnine} in Python. & R, Python \\
004 Analyze data & Produces a sample-selection table, descriptive
statistics, a correlation matrix, and a main regression table with fixed
effects, exported as \LaTeX{}, Word, and RTF. & R, Python, Stata \\
005 Data provenance & Exports the sample-identifier file in Parquet and
CSV; computes SHA256 hashes for every raw, derived, and output file. & R,
Python \\
\bottomrule
\end{tabularx}
\end{table}

An orchestrator script chains the steps via a \texttt{batch\_run()}
helper that spawns each script in a fresh child process and writes a
uniform plain-text log per script. Running each step in a fresh process
is itself a teaching point: it guarantees that no step depends on state
left in memory by an earlier one, a failure mode that interactive
notebook- and console-centered work habits make chronically likely. When
the user has selected a Stata-inclusive language combination, the
orchestrator additionally invokes the Stata analysis script through a
helper that produces a log of the same shape, so the artifact set is
uniform regardless of language mix.

\subsection{Polyglot design and cross-language consistency}

The template's most unusual design decision is shipping parallel
implementations of every step in multiple languages. Each numbered step
exists as a matched pair of R and Python scripts (with Stata additionally
covering the analysis step), written to be structurally parallel: the
same step boundaries, the same intermediate files, the same variable
names, and the same outputs.

This design serves three pedagogical purposes. First, it removes the
language-choice barrier to adoption: a student or instructor can work
entirely in their incumbent language while gaining the workflow
conventions, rather than facing a bundled language conversion on top of a
workflow conversion. Second, the matched scripts function as a Rosetta
stone. A student who knows Stata can read the R implementation of a
regression they already understand---\texttt{reghdfe} with
\texttt{estout} \citep{reghdfe, estout} beside \texttt{fixest} with
\texttt{modelsummary} \citep{fixest, modelsummary} beside
\texttt{pyfixest} \citep{pyfixest}---and transfer their conceptual model
across syntax. In our experience this is among the fastest ways for a
student to acquire a second research language.

Third, and most subtly, the parallel implementations teach a verification
habit. The templates demonstrate that the same regression produces the
same coefficients, standard errors, and sample counts in all three
languages, and the materials treat any cross-language discrepancy in
observation counts as a bug to be traced rather than noise to be
tolerated. Students internalize sample-count reconciliation as a routine
check---a professionally important habit in a field where subtle join and
filter differences across implementations are a common source of
irreproducibility.

The data layer supports the polyglot design: intermediate datasets are
stored as parquet files, which R, Python, and (via export) Stata read
identically, avoiding the lossy CSV round-trips that otherwise accumulate
at language boundaries.

\subsection{Journal-policy artifacts as a by-product}

A central design principle is that policy compliance should fall out of
running the pipeline rather than being a separate documentation chore
performed at submission time---the point at which, in practice, it is
hardest to reconstruct. Each element of the \textit{Journal of Accounting
Research} artifact triad \citep{JARpolicy} maps to a mechanism in the
template:

\begin{enumerate}[leftmargin=*]
  \item \textbf{Code from raw inputs.} The numbered pipeline \emph{is}
  the required code; because step 001 downloads directly from WRDS, the
  pipeline's raw inputs are the licensed databases themselves and the
  reproduction chain has no undocumented head.
  \item \textbf{Execution logs.} The orchestration layer writes one
  uniform plain-text log per script, in the same format across R, Python,
  and Stata, capturing the end-to-end execution record the policy
  requires.
  \item \textbf{Sample identifiers.} Step 005 exports the identifying
  keys (firm and period identifiers) of every observation in the final
  regression sample, in both Parquet and CSV.
\end{enumerate}

Step 005 additionally computes a SHA256 hash inventory of every raw,
derived, and output file. The hash inventory goes beyond current policy
requirements and gives students a concrete grip on an otherwise abstract
idea: two research states are identical exactly when their hashes match,
so ``did anything change?'' becomes a checkable question rather than a
feeling.

\section{Learning objectives and instructional design}
\label{sec:design}

After working through the hub and completing the pipeline in at least one
language, learners should be able to:

\begin{itemize}[leftmargin=*]
  \item Structure an empirical research project across a version-control
  repository, a local working clone, a cloud-synchronized data directory,
  and a \LaTeX{} collaboration platform, with code, data, and manuscript
  artifacts cleanly separated.
  \item Use Git and GitHub for everyday research workflows: committing
  meaningful changes, branching for alternative specifications, tagging
  paper versions for journal submissions, and collaborating with
  coauthors via pull requests.
  \item Configure a reproducible development environment, including
  project-level configuration via a \texttt{.env} file, Python
  environments via \texttt{uv} \citep{uv}, and operating-system keyring
  storage of database credentials.
  \item Implement a five-step empirical pipeline---licensed-database
  download, transformation, figures, analysis tables, and provenance
  reporting---in R, Python, Stata, or a working combination.
  \item Produce publication-grade tables in \LaTeX{}, Word, and RTF
  formats from the same analysis code, and integrate the \LaTeX{} outputs
  into a journal-ready manuscript.
  \item Produce the artifact set required by journals' data and code
  sharing policies: per-script execution logs, sample-identifier files,
  and file-hash inventories.
  \item Apply analogous conventions in SAS where a project requires it,
  using the reference macro library in the hub.
\end{itemize}

Several instructional design decisions deserve comment because they
depart from common practice in workflow teaching.

\paragraph{Backward design from the journal artifact set.}
Rather than teaching reproducibility as a general virtue and hoping
students map it onto their field's requirements, the materials treat the
journal policy artifacts as the terminal assessment and derive the
curriculum from them. This choice exploits an unusually clean alignment
available in this field: the professional gatekeeping mechanism (journal
policy) already specifies concrete, checkable outputs, which gives
workflow instruction the kind of authentic assessment target that the
education literature generally has to construct artificially.

\paragraph{Scripts and IDEs rather than notebooks.}
The materials deliberately center scripted pipelines run from an IDE
rather than notebooks, because the terminal artifact---a
policy-compliant, batch-executable research pipeline---is what students
must ultimately produce, and because fresh-process batch execution
surfaces the hidden-state errors that interactive work styles conceal.
The IDE chapter treats editor choice as a configuration decision rather
than an identity, reflecting the multi-language design.

\paragraph{Explicit treatment of configuration and secrets.}
Hardcoded absolute paths and database passwords in scripts are the two
most common workflow defects in student (and faculty) research code. The
materials confront both directly: every script resolves data locations
through a project-level \texttt{.env} file read identically by all four
languages, and WRDS credentials live in the operating-system keyring,
never in code. These conventions also make the difference between ``runs
on my machine'' and ``runs on my coauthor's machine'' concrete on day
one.

\paragraph{AI assistants as a first-class workflow topic.}
Each repository carries a context file documenting project conventions
for AI coding assistants, and the hub includes a chapter on the emerging
cross-tool standard for such files. Doctoral students are already using
AI assistants for research code; the materials take the position that
assistant-readable project documentation is now part of a reproducible
project's infrastructure, and that the discipline of writing down project
conventions benefits human collaborators regardless.

\section{Classroom use}
\label{sec:classroom}

The materials were developed for, and have been iteratively refined
through, doctoral teaching in the accounting program at \institution,
where earlier versions were used in a PhD seminar on empirical research
methods.
%%%%%%%%%%%%%%%%%%%%%%%%%%%%%%%%%%%%%%%%%%%%%%%%%%%%%%%%%%%%%%%%%%%%%%%%%%%%
% AUTHOR TO COMPLETE: number of semesters and course numbers, approximate
% student counts, and any student feedback you are comfortable
% characterizing. Reviewers at an education journal will look for this.
%%%%%%%%%%%%%%%%%%%%%%%%%%%%%%%%%%%%%%%%%%%%%%%%%%%%%%%%%%%%%%%%%%%%%%%%%%%%
\textbf{[AUTHOR: semesters taught, course numbers, approximate
enrollment, and characterization of student feedback.]}
Many of the underlying examples and conventions were first developed in
published empirical work in the field
\citep{weisbrod2019stockholders, bochkay2022roles, larocqueconsensus},
and the public companion repository for \citet{larocqueconsensus} is
referenced throughout the hub as a real-world, full-scale instance of the
conventions the templates teach in miniature.

The recommended course sequence assigns the hub's topic chapters as
preparatory reading, then has each student create a personal repository
from the pipeline template, select a language combination at setup, and
run the pipeline end to end against WRDS. Because step 001 downloads
directly from WRDS, students must have institutional WRDS
credentials---standard at research universities in business
fields---and the first class session doubles as a credential and
environment setup clinic. A capstone exercise extends the pipeline with
an alternative specification developed on a Git branch and produces a
journal-style manuscript draft from the students' own table and figure
outputs using the manuscript template.

The materials also support two lighter-weight adoption modes observed in
practice: independent doctoral study, in which a student works through
the hub and template alone as preparation for a first research project;
and lab-onboarding use, in which an advisor points incoming coauthors at
specific topic chapters (most commonly the Git and project-structure
chapters) as shared vocabulary for collaboration conventions.

\section{Discussion}
\label{sec:discussion}

\subsection{Adoptability beyond accounting and finance}

Although the running example and the journal-policy framing are specific
to accounting and finance, most of the materials transfer directly to
adjacent observational-data fields---economics, political science,
empirical legal studies---where the same stack (licensed or restricted
administrative data, multi-language teams, code-and-log based
verification) increasingly prevails. An instructor outside business
fields would need to replace the WRDS download step with their field's
data source and would likely retarget the journal-policy discussion at
their own field's replication norms
\citep{christensen2018transparency}. The topic chapters (project
structure, Git, IDEs, environments, configuration) are field-agnostic as
written. The CC-BY-4.0 license permits adaptation and redistribution with
attribution, and the template mechanism means adopting instructors can
maintain institution-specific variants without forking histories.

\subsection{Teaching reproducibility when the data cannot be shared}

The materials offer the statistics and data science education community a
worked example of a problem it has largely set aside: what reproducibility
instruction looks like when the data are licensed and cannot be
published. The literature's canonical practices---publish the data with
the analysis, make the repository self-contained, let anyone press
``run'' \citep{peng2011reproducible, dogucu2022reproducible}---assume
shareable data, and even guidance aimed at sensitive data addresses
material that can be deidentified and then released
\citep{horton2024editorial}. Commercially licensed data admit neither
remedy: the vendor's contents cannot be deidentified into shareability.

Under licensed data, verifiability must instead rest on a deterministic
pipeline anchored at the licensed source, execution logs, sample
identifiers, and content addressing; that is, on \emph{evidence about the
computation} rather than the ability of arbitrary readers to re-execute
it. This artifact-based model of reproducibility is not a second-best
workaround but the institutionalized standard in a major research field,
and teaching it requires its own materials. Instructors whose students
will work with restricted data---health records, administrative
microdata, commercial licenses---may find the pattern, and the concrete
artifact set, directly reusable.

\subsection{Limitations and future development}

The materials have known limitations. First, systematic evidence of
learning outcomes is limited to instructor observation and student
feedback; the classroom claims here are an experience report rather than
an efficacy study. Second, the pipeline demonstrates a single research
design; students whose first project departs far from an archival event
study must adapt more of the skeleton themselves, although the workflow
layer is design-agnostic. Third, the SAS support is deliberately a
reference macro library rather than a fourth parallel pipeline,
reflecting the language's declining share of new projects but continuing
presence in legacy code. Planned development includes additional worked
designs and expanded topic chapters; because the hub is versioned and
archived with a concept DOI, adopters can pin or track versions as they
prefer.

\section{Conclusion}

The suite described here packages the workflow layer of empirical
accounting and finance research---project structure, version control,
environments and configuration, a scripted multi-language pipeline, and
journal-policy artifact production---into an open, adoptable teaching
resource built around an authentic research design. It extends the
statistics and data science education community's reproducibility agenda
into business doctoral education, and contributes back a worked model of
reproducibility instruction under licensed data.

\ifanon
\else
\section*{Acknowledgements}

The author thanks doctoral students and colleagues who have used and
provided feedback on earlier versions of these materials.
\fi

\section*{Data availability statement}

No original data were collected for this article. The teaching materials
are openly available under a CC-BY-4.0 license
\ifanon
at repositories whose URLs are withheld for anonymous review and will be
supplied to the editors on request; the hub is archived with a citable
concept DOI.
\else
at \hubsite, with companion templates at \pipelinerepo{} and
\manuscriptrepo. The hub is archived at Zenodo under concept DOI
\href{https://doi.org/10.5281/zenodo.20499216}{10.5281/zenodo.20499216}.
\fi
The running example uses licensed data from Wharton Research Data
Services (WRDS), available to researchers at subscribing institutions.

\section*{Disclosure statement}

The author reports no conflict of interest.

\newpage
\bibliography{paper}

\end{document}
